实偏导可交换，故有 Hermitian 对称性。对角元中，$\alpha\ne\beta$ 对应的项两两抵消，$\alpha=\beta$ 对应的每一项的系数均为一。由此得到迹公式。将 Hessian 映射作用于 $\xi\xi^{\mathsf t}$，得到 $gg^*$。任意非负实型均可写成这种形式之和。通常的乘积法则给出 $H(uv)$ 的恒等式。在四槽情形下，四元数酉矩阵 $U$ 的实转置对编码余向量的作用是 $U^*$。秩一情形的计算给出 $H(u\circ U)=U^*(H(u)\circ U)U$，再由线性性推广到任意实 Hessian。作用于坐标指标的实正交矩阵保持三槽子空间，因此也得到三槽情形下的协变性。

记 $P_a=p_a+\mathbf k q_a$，故 $2\partial_{\bar P_a}=\partial_{p_a}+\mathbf k\partial_{q_a}$。因为 $\mathbf{ki}=\mathbf j$，梯度为 $2(\partial_{\bar P_a}u)\mathbf i$。在 Hessian 中，与导数 $(p_ap_b,p_aq_b,q_ap_b,q_aq_b)$ 对应的四个系数为 $1,-\mathbf k,\mathbf k,1$，得到 $4\partial_{\bar P_a}\partial_{P_b}u$。对全纯映射 $F$，复链式法则给出
\[
\partial_{\bar P_a}\partial_{P_b}\sum_\nu|F_\nu|^2
 =\sum_\nu\overline{J_{\nu a}}\,J_{\nu b},\qquad
G(\rho\circ F)=2J^*(\partial_{\bar F_\nu}\rho)_\nu\,\mathbf i,
\]
其中 $\rho$ 的各导数在 $F(P)$ 处取值。
